\documentclass[11pt,a4paper]{article}
\usepackage[T1]{fontenc}
\usepackage[margin=2.5cm]{geometry}
\usepackage{xcolor}
\usepackage{amsmath}
\usepackage{booktabs}
\usepackage{array}
\usepackage{adjustbox}
\usepackage{placeins}
\usepackage{eso-pic}
\usepackage{pgfplots}
\usepgfplotslibrary{groupplots}
\pgfplotsset{compat=1.17}
\usepackage[numbers,sort&compress]{natbib}
\usepackage[hidelinks]{hyperref}
\hypersetup{pdfauthor={}, pdftitle={Supplementary material}}
% Build switches. The build script sets them in the file buildconfig.tex.
\newif\ifsubmission \submissionfalse
\newif\ifresultsloaded \resultsloadedfalse
\newif\ifdryrun \dryrunfalse
\newif\ifposthoc \posthocfalse
\providecommand{\resultsdir}{results}
\InputIfFileExists{buildconfig}{}{}

% An unfinished passage. The draft shows it in a red box. The submission build stops at it.
\newcounter{pendingcount}
\newcommand{\pending}[1]{%
  \ifsubmission
    \GenericError{}{Submission build: an unfinished passage remains}{Remove every unfinished passage.}{#1}%
  \else
    \stepcounter{pendingcount}%
    \par\smallskip\noindent
    {\color{red}\fbox{\parbox{\dimexpr\linewidth-2\fboxsep-2\fboxrule\relax}{\small\textbf{PENDING:} #1}}}\par\smallskip
  \fi}

% The mark of a table or figure that comes from synthetic test input.
\newcommand{\dryrunmark}{\textcolor{red}{\textbf{DRY RUN. SYNTHETIC TEST INPUT. NOT A RESULT.}}}

% A generated result file. A missing file gives a visible note in the draft and stops the submission build.
\newcommand{\resultfile}[1]{%
  \InputIfFileExists{\resultsdir/#1}{}{%
    \ifsubmission
      \GenericError{}{Submission build: the result file #1 is missing}{Run the results script.}{}%
    \else
      \pending{The generated result \texttt{\detokenize{#1}} is not available. The results script writes it from the bench analysis.}%
    \fi}}
% The same for a result that is optional in the draft.
\newcommand{\optionalresultfile}[1]{\InputIfFileExists{\resultsdir/#1}{}{}}

\newcommand{\loadresultnumbers}{%
  \InputIfFileExists{\resultsdir/numbers}{}{}%
  \ifsubmission
    \ifresultsloaded\else
      \GenericError{}{Submission build: no results are loaded}{Run the results script.}{}%
    \fi
    \ifdryrun
      \GenericError{}{Submission build: the results are a dry run}{Use the measured results.}{}%
    \fi
  \fi}

% A page mark for every draft page that holds dry-run content.
\newcommand{\dryrunwatermark}{%
  \ifdryrun
    \AddToShipoutPictureFG{\AtPageUpperLeft{\raisebox{-1.2cm}{\hspace{1.5cm}\dryrunmark}}}%
  \fi}

\loadresultnumbers
\dryrunwatermark
\renewcommand{\thesection}{S\arabic{section}}
\renewcommand{\thetable}{S\arabic{table}}
\renewcommand{\thefigure}{S\arabic{figure}}
\renewcommand{\theequation}{S\arabic{equation}}
\title{Supplementary material\\
\large A camera-referenced robotic bench method for comparing pen digitisers,\\ applied to an Apple Pencil and a Wacom tablet}
\author{Anonymous authors}
\date{}
\begin{document}
\maketitle

\section{Movement programs and fault programs}
\label{sec:s-programs}
Coordinates are in millimetres. The origin is at the top left, $x$ goes to the right, and $y$ goes down.
The loops curve has the parameter $t$ from 0 to $16\pi$:
\begin{equation}
x(t)=8+\frac{12\,t}{2\pi}+4\sin t,\qquad y(t)=24-10\,(1-\cos t).
\end{equation}
The spiral curve has the angle $\theta$ from 0 to $4\pi$ and the radius $r(\theta)=8+21\,\theta/(2\pi)$:
\begin{equation}
x(\theta)=52+r\cos\theta,\qquad y(\theta)=52-r\sin\theta.
\end{equation}
A program at scale $a$ multiplies both coordinates by $a$.
The line runs from $(4,8)$ to $(128,8)$.
Its cruise covers 120~mm in 6~s.
The position along the line carries the added term $A\,(1-\cos 2\pi f t)$, so the speed is $20+2\pi f A\sin 2\pi f t$ in mm/s.
The cruise holds a whole number of cycles.

On the loops, the speed is proportional to the cube root of the radius of curvature, with the radius limited to 40~mm.
The constant of proportionality makes the mean speed outside the ramps equal to the program speed.
The commanded speed of the loops programs ranges from 6.8 to 26.9~mm/s at 20~mm/s, and up to 107.7~mm/s at 80~mm/s.

A program is sampled every 10~ms into timed moves.
The moves follow the oscillation of a line program with straight segments.
The band power of the three line programs on the move list is therefore 1.2, 3.1, and 5.8\,\% below the analytic value of the 5, 8, and 11~Hz programs.
The pilot rule of the main text uses the move-list values as the command.
The features resample every record to a 200~Hz grid by linear interpolation.
For a noise-free record at 240 samples per second, this lowers the band power by a further 0.7 to 3.1\,\%.
This effect is part of the measurement chain under test.

\paragraph{Fault programs.}
The clean median of a block is the median over its valid runs of the loops program at scale 1.0 and 20~mm/s.
A fault run counts as detected when every criterion for its device holds.
\begin{itemize}
\item Dwell fault. The motion duration is within 2\,\% of the clean median plus 0.5~s. On device W, the recording has no gap in the contact samples, because the tablet reports a still pen. On device P, the expected result (a gap with two contact segments, or no gap) is written from the pilot before the lock.
\item Lift fault. The recording has two contact strokes.
\item Truncation fault. The path length and the motion duration are at most 0.6 times their clean medians.
\end{itemize}

\paragraph{Quality checks of a recording.}
A contact stroke is split at a gap above 0.05~s.
A segment shorter than 0.1~s is dropped.
A recording with a contact sampling rate below 40~Hz is rejected.

\section{Camera procedures}
\label{sec:s-camera}
\paragraph{Requirements.}
The camera captures 240 frames per second without interpolated frames.
The field of view covers $152\times124$~mm in the marker plane.
The optical axis is normal to the device surface within $2^\circ$.
The lighting has no flicker.
Focus, exposure, and white balance are locked or held constant by the procedure.

\paragraph{Frame clock.}
The LED is recorded for 60~s once before the lock.
The frame count per LED second gives the frame rate.
In every clip, each LED period inside the motion interval must hold the frames of the measured rate within one frame plus 0.1\,\%, and all periods together must hold their total within the same tolerance.
A frame that is identical to the previous frame is a duplicate.
A clip with a duplicated or missing frame inside the motion interval is not a valid camera record.
Where the file holds capture timestamps, every frame interval inside the motion interval must be within 1\,\% of the median interval plus 1~\textmu s.

\paragraph{Rolling shutter.}
A rolling shutter reads the image rows one after another.
A 1~kHz LED then appears as bands, and the band pitch in rows gives the row readout time.
The movement of the bands between frames gives the readout direction.
The time of each marker position is its frame time plus its row index times the row readout time, counted in the readout direction.
In a test with rendered frames of the fast loops program, the RMS error of the track was 0.004~mm with the correction, 0.013~mm without it, and 0.026~mm with the wrong direction.

\paragraph{Tracking.}
Until the tracker first finds the marker, it searches a square of 24~mm around the commanded start point.
After that, each frame searches a square of 10~mm around the last position.
The marker is the largest bright region that lies completely inside the square, has 0.5 to 2 times the area of the marker dot, and is surrounded by the dark field.
The centroid is corrected for lens distortion, mapped to the marker plane, and shifted by the marker offset.

\paragraph{Alignment of the camera track.}
A coarse alignment slides the commanded speed profile over the clip and takes the lag with the largest correlation.
Both profiles are smoothed with a moving average of 0.2~s for this step only.
The motion interval of the camera track is then searched inside the commanded motion interval at this lag, widened by 0.2~s before and 0.15~s after.
The camera reference therefore never depends on the device under test.
A camera motion interval that touches an edge of this window is a motion failure.
A coarse correlation below 0.5 is a motion failure when the clip passes every camera check.

For the spatial comparison, the camera track is aligned to the device record at the lag, up to 0.5~s, that maximises the correlation of the two speed profiles, refined below one frame by a parabolic fit.

\paragraph{Fit to the commanded path.}
Each sample is first paired with the commanded position at its commanded time, which gives a first rigid fit.
An iterative closest-point fit to the commanded path then repeats the rigid fit with nearest points~\cite{Besl1992}.

\section{Precision of the design}
\label{sec:s-precision}
Tables~\ref{tab:s-program} to \ref{tab:s-small} give the simulation results that the main text summarises.
A run value is a device mean plus a mounting effect plus run noise.
$\sigma_\mathrm{mount}$ and $\sigma_\mathrm{run}$ are the standard deviations of the two terms.
All values are in units of the limit $L$.
The mounting effect is either shared by all programs of a mounting or drawn for each program.

\begin{table}[ht]
\centering
\caption{Per-program decision with six mountings and five runs: share of 20\,000 simulated studies in which the program is declared equivalent at a level of 0.05, for a true bias of 0, $0.5L$, and $L$.}
\label{tab:s-program}
\begin{tabular}{@{}rrrrrr@{}}
\toprule
$\sigma_\mathrm{run}$ & $\sigma_\mathrm{mount}$ & Median half-width & Bias 0 & Bias $0.5L$ & Bias $L$ \\
\midrule
0.25 & 0.10 & 0.155 & 1.000 & 1.000 & 0.048 \\
0.25 & 0.25 & 0.282 & 1.000 & 0.896 & 0.049 \\
0.25 & 0.50 & 0.527 & 0.855 & 0.456 & 0.046 \\
0.50 & 0.25 & 0.346 & 0.998 & 0.763 & 0.045 \\
0.50 & 0.50 & 0.564 & 0.797 & 0.421 & 0.049 \\
\bottomrule
\end{tabular}
\end{table}

\begin{table}[ht]
\centering
\caption{Feature decision: share of simulated studies in which every program of the set is declared equivalent. Each entry gives the share at a level of 0.05 and at 0.05/6, the strictest Holm level. The bias applies to every program.}
\label{tab:s-feature}
\begin{tabular}{@{}rrrllll@{}}
\toprule
$\sigma_\mathrm{run}$ & $\sigma_\mathrm{mount}$ & Bias & Mounting effect & 12 programs & 5 loops programs & 3 line programs \\
\midrule
\multicolumn{7}{@{}l}{Six mountings, five runs} \\
0.25 & 0.25 & 0 & shared & 1.000 / 0.973 & 1.000 / 0.982 & 1.000 / 0.988 \\
0.25 & 0.25 & 0 & per program & 0.999 / 0.940 & 1.000 / 0.975 & 1.000 / 0.986 \\
0.25 & 0.25 & $0.5L$ & shared & 0.727 / 0.287 & 0.790 / 0.371 & 0.828 / 0.436 \\
0.25 & 0.25 & $0.5L$ & per program & 0.271 / 0.002 & 0.578 / 0.075 & 0.719 / 0.205 \\
0.25 & 0.50 & 0 & shared & 0.727 / 0.183 & 0.769 / 0.237 & 0.791 / 0.272 \\
0.25 & 0.50 & 0 & per program & 0.158 / 0.000 & 0.459 / 0.008 & 0.632 / 0.058 \\
0.50 & 0.25 & 0 & shared & 0.981 / 0.625 & 0.992 / 0.780 & 0.995 / 0.848 \\
0.50 & 0.25 & 0 & per program & 0.978 / 0.472 & 0.991 / 0.735 & 0.994 / 0.826 \\
0.50 & 0.25 & $0.5L$ & shared & 0.317 / 0.035 & 0.466 / 0.094 & 0.562 / 0.158 \\
0.50 & 0.25 & $0.5L$ & per program & 0.041 / 0.000 & 0.266 / 0.011 & 0.450 / 0.065 \\
\addlinespace
\multicolumn{7}{@{}l}{Eight mountings, five runs} \\
0.25 & 0.25 & $0.5L$ & shared & 0.883 / 0.545 & 0.917 / 0.635 & 0.934 / 0.683 \\
0.25 & 0.25 & $0.5L$ & per program & 0.650 / 0.063 & 0.834 / 0.321 & 0.899 / 0.504 \\
0.25 & 0.50 & 0 & shared & 0.908 / 0.516 & 0.927 / 0.577 & 0.937 / 0.618 \\
0.25 & 0.50 & 0 & per program & 0.609 / 0.018 & 0.813 / 0.184 & 0.880 / 0.360 \\
\bottomrule
\end{tabular}
\end{table}

\begin{table}[ht]
\centering
\caption{Correction with covariates on the 12 programs, with a true intercept of $0.5L$, a true slope of 1.05, six mountings, five runs, and $\sigma_\mathrm{run}=0.25L$. The \emph{linear} truth has no program effect. The \emph{speed} truth adds $0.5L$ for each doubling of the speed, which the model can represent. The \emph{program} truth adds offsets of $\pm0.5L$ that alternate between programs, which the model cannot represent.}
\label{tab:s-correction}
\begin{tabular}{@{}rlrrrr@{}}
\toprule
$\sigma_\mathrm{mount}$ & Truth & Largest & Check passes & Largest true residual & Jackknife SE \\
 & & uncorrected bias & on all programs & bias of the fit (median) & of $\beta$ (median) \\
\midrule
0.10 & linear & 1.05 & 1.000 & 0.093 & 0.0265 \\
0.10 & speed & 2.05 & 1.000 & 0.093 & 0.0267 \\
0.10 & program & 1.55 & 0.263 & 0.849 & 0.0553 \\
0.25 & linear & 1.05 & 0.999 & 0.154 & 0.0326 \\
0.25 & speed & 2.05 & 0.999 & 0.156 & 0.0331 \\
0.25 & program & 1.55 & 0.033 & 0.895 & 0.0823 \\
0.50 & linear & 1.05 & 0.726 & 0.279 & 0.0795 \\
0.50 & speed & 2.05 & 0.721 & 0.276 & 0.0799 \\
0.50 & program & 1.55 & 0.003 & 1.016 & 0.1344 \\
\bottomrule
\end{tabular}
\end{table}

On the 12 programs, with a device difference that the model can represent, the lack-of-fit check passes on every program in 99.9\,\% of studies for $\sigma_\mathrm{mount}$ up to $0.25L$.
With program offsets that the model cannot represent, it passes in 26\,\% of studies or fewer.
On the 12 programs, the check therefore detects a misfit of this size.
It does not estimate the bias of the correction on new data.

Table~\ref{tab:s-small} gives the same simulation for the two smaller program sets, with 5000 simulated studies per row.
On these sets, the check passes program offsets of $\pm0.5L$ in about 96\,\% of studies at $\sigma_\mathrm{mount}=0.1L$.
These passes are correct, because the true residual bias of about $0.67L$ lies inside the limit.
The check detects such offsets only at $\sigma_\mathrm{mount}$ of $0.25L$ or more.
On the loops programs, the slope takes up a dependence on speed: with the speed truth, the fitted slope is 1.142 against a true slope of 1.05.

\begin{table}[ht]
\centering
\caption{Correction on the two smaller program sets: the slope model on the five loops programs and the shift model on the three line programs. The truths are those of Table~\ref{tab:s-correction}. The line programs have one speed, so their speed truth equals the linear truth.}
\label{tab:s-small}
\begin{tabular}{@{}lrlrrrl@{}}
\toprule
Set and model & $\sigma_\mathrm{mount}$ & Truth & Largest & Check passes & Largest true & Median $\beta$ \\
 & & & uncorrected bias & on all programs & residual bias (median) & \\
\midrule
Loops, slope & 0.10 & linear & 1.05 & 1.000 & 0.079 & 1.050 \\
Loops, slope & 0.10 & speed & 2.05 & 1.000 & 0.172 & 1.142 \\
Loops, slope & 0.10 & program & 1.55 & 0.960 & 0.673 & 1.095 \\
Loops, slope & 0.25 & linear & 1.05 & 1.000 & 0.135 & 1.050 \\
Loops, slope & 0.25 & speed & 2.05 & 0.999 & 0.230 & 1.142 \\
Loops, slope & 0.25 & program & 1.55 & 0.384 & 0.730 & 1.095 \\
Loops, slope & 0.50 & linear & 1.05 & 0.806 & 0.239 & 1.050 \\
Loops, slope & 0.50 & speed & 2.05 & 0.662 & 0.340 & 1.142 \\
Loops, slope & 0.50 & program & 1.55 & 0.041 & 0.829 & 1.094 \\
Line, shift & 0.10 & linear & 0.45 & 1.000 & 0.046 & fixed \\
Line, shift & 0.10 & program & 0.95 & 0.958 & 0.668 & fixed \\
Line, shift & 0.25 & linear & 0.45 & 1.000 & 0.101 & fixed \\
Line, shift & 0.25 & program & 0.95 & 0.439 & 0.670 & fixed \\
Line, shift & 0.50 & linear & 0.45 & 0.850 & 0.199 & fixed \\
Line, shift & 0.50 & program & 0.95 & 0.051 & 0.732 & fixed \\
\bottomrule
\end{tabular}
\end{table}

These simulations use an assumed spread of the device W values between programs, about $13L$ on the loops programs.
The designed spread is smaller for the speed CV.
At the designed spread of the speed CV, the median jackknife standard error of $\beta$ is 0.054 at $\sigma_\mathrm{mount}=0.1L$ and 0.064 at $0.25L$.
At the designed spread of the loop height, it is 0.0075.
The limit of $0.25L$ for $\sigma_\mathrm{mount}$ is 0.9 to 1.2\,\% for the four camera-referenced features, 0.0075 for the speed CV, and 6.7\,\% for the band power.

\section{Consequence experiment: further settings}
\label{sec:s-consequence}
\paragraph{Generator.}
For participant $i$, session $j$, and a feature with limit $L$, let $s=L/(\sqrt{2}\,\Phi^{-1}(0.55))$. Then
\begin{equation}
Y_i\sim\mathrm{Bernoulli}(0.30),\qquad
X_{ij}=\mu+s\left(\delta\,(Y_i-0.30)+\sqrt{0.80}\,U_i+\sqrt{0.20}\,V_{ij}\right),
\end{equation}
where $U_i$ and $V_{ij}$ are independent standard normal variables, $\delta\in\{0,0.8\}$ is the effect, and $\mu$ is the noise-free feature value of the selected bench program.
One random draw of the session error is shared by the two devices and scaled by the run standard deviation of each device.
The mounting effects are drawn independently for the two devices.

\paragraph{Assumed grid.}
The noise settings give the standard deviations between runs and between mountings in units of $L$.
\emph{Zero}: both devices $(0,0)$.
\emph{Nominal}: P $(0.25,0.10)$ and W $(0.25,0.25)$.
\emph{Stress}: P $(0.25,0.25)$ and W $(0.50,0.50)$.
In the grid, the correction is the known shift $b$.
It removes the bias and leaves the noise.
A control with a wrong correction (an error of $0.5L$ at the centre, with a slope of 1.05 where the model has a slope) runs at $k\in\{-1,0,1\}$ with nominal noise.

\paragraph{Program selection.}
For each feature, the bench program is chosen before any label or model result is read.
The geometry must match the target.
Among the matching programs, the choice minimises $|\ln(a^*/a)|+|\ln(q^*/q)|$ for the target scale $a^*$ and speed $q^*$.
For the line programs, the choice first minimises the difference of the oscillation frequency.

\paragraph{Validation and metrics.}
The AUC compares only cases and controls of one test fold, and the fold values are averaged with weights by pair count.
In record-wise validation, each participant has a total weight of one.
The calibration intercept is the intercept of a logistic regression of the label with the log-odds of the prediction as an offset.
The calibration slope comes from a regression with a free intercept.
The net benefit at a threshold $t$ is the share of true positives minus the share of false positives times $t/(1-t)$.
A material change is an absolute difference from the reference of at least 0.05 for the AUC, 0.20 for the calibration intercept or slope, and 0.01 for the mean net benefit.
These thresholds are descriptive judgements.
A setting with fewer than 95\,\% usable cohort pairs is marked incomplete.

\paragraph{Sensitivity analyses.}
One setting changes at a time: a cohort size of 200 or 1200, a case share of 0.08, an effect of 0.4, a between-person variance share of 0.95, a nearest-neighbour model, a model with all six features, and device assignment that is independent of or identical to case status.
With identical assignment, device-held-out validation has only one class in each test set, and its metrics are reported as unavailable.
For the speed CV, generated values below zero are kept, and their share is reported; one analysis sets them to zero.
With the measured bench inputs, two analyses use the lower and the upper bound of each bias.

\section{Additional results}
\label{sec:s-results}
\optionalresultfile{tab_bland_altman}
\optionalresultfile{fig_bland_altman}
\optionalresultfile{tab_selection}
\optionalresultfile{tab_consequence_grid}
\ifresultsloaded\else
\pending{The additional result tables and figures are not available. The results script writes them from the bench analysis and the consequence experiment.}
\fi

\bibliographystyle{unsrtnat}
\bibliography{references}
\end{document}
